\chapter{项目概述}

\section{项目背景}

本项目是“操作系统设计综合实践”课程的小组课题，开发基线为 xv6-RISC-V。xv6 原有调度器采用轮转方式遍历进程表，对所有可运行进程一视同仁。这种实现结构清楚，适合教学，但无法表达任务之间的优先级差异。项目在保留 xv6 进程状态转换和锁机制的前提下，引入静态优先级调度，并使用 Aging 缓解低优先级进程长期等待的问题。

项目规定优先级取值范围为 0 至 31，数值越小表示优先级越高，默认值为 10。实现范围包括进程优先级字段、fork 继承、静态优先级调度、同级轮转、Aging、优先级系统调用和用户命令。
代码仓库地址为：\href{https://github.com/Aphrodite-deepblue/os-design-practice}{\texttt{github.com/Aphrodite-deepblue/os-design-practice}}。

\section{小组分工}

小组按内核数据、调度策略、Aging、系统调用和测试五个模块分工。每项任务都对应明确的代码区域和 Git 提交，分工见表\ref{tab:team-work}。

\begin{table}[htbp]
  \centering
  \small
  \caption{第 35 组成员及分工}
  \label{tab:team-work}
  \begin{tabularx}{\textwidth}{C{1.5cm}C{2.4cm}C{2.5cm}X}
    \toprule
    姓名 & 学号 & GitHub 账号 & 主要任务 \\
    \midrule
    赖弈铭 & 3022002026 & \href{https://github.com/Aphrodite-deepblue}{\texttt{Aphrodite-deepblue}} & Aging、等待计数与防饥饿逻辑 \\
    刘文睿 & 3023001172 & \href{https://github.com/milabo0718}{\texttt{milabo0718}} & 静态优先级调度与同优先级轮转 \\
    吴雨桐 & 3022001567 & \href{https://github.com/Curtai123}{\texttt{Curtai123}} & 测试设计、集成验证与结果整理 \\
    党正清 & 3023244304 & \href{https://github.com/dzq180}{\texttt{dzq180}} & 优先级系统调用与 \texttt{nice} 命令 \\
    袁晟棋 & 3023244399 & \href{https://github.com/yuan-sq}{\texttt{yuan-sq}} & 进程优先级数据、fork 继承与集成 \\
    \bottomrule
  \end{tabularx}
\end{table}

\section{AI 使用情况}

项目开发和材料整理过程中使用了 OpenAI Codex（GPT-5 系列模型）。AI 的使用范围是读取课程要求与仓库现状、辅助定位 xv6 系统调用接入位置、复核边界条件和锁语义、整理本地与 CI 测试命令，以及起草和校正文档。AI 没有替代小组确定调度语义，也没有生成虚构的终端输出、截图或实验指标。

成员 4 在实现优先级控制接口时，使用 Codex 辅助核对从系统调用号、内核分发表、参数入口、进程表操作到用户态调用桩的完整链路，并检查 \texttt{nice} 的数字解析、PID 校验和错误处理。最终采用的接口仍以小组统一约定为准：优先级范围为 0 至 31，设置成功后清零 \texttt{wait\_ticks}，进程字段访问受 \texttt{p->lock} 保护。相关实现由成员逐项复核，并形成提交 \texttt{43b682f} 和 \texttt{30603f0}。

AI 给出的建议不能作为功能正确性的证据。成员 4 的修改实际完成了交叉编译、QEMU 启动、\texttt{nice} 查询与设置、合法边界和非法输入检查，并执行了原有 \texttt{usertests} 与 \texttt{crash} 回归。后续测试模块又在真实 QEMU 中运行 \texttt{prioritytest}、\texttt{agingtest} 和 \texttt{usertests -q}，运行器保存原始日志、结构化结果和源码哈希；报告只引用这些真实结果，不将 AI 推断写成实验结论。

在测试与汇报阶段，Codex 还用于核查公开代码、CI 配置和文档一致性，并辅助草拟说明文字。所有建议都由成员结合当前 \texttt{main} 分支复核，最终 Git 作者、提交内容、实验执行和报告结论由小组成员负责。必要的使用场景和人工验证过程以本节、开发记录及 Git 历史为准，不用大段原始对话替代可复核的工程证据。
